\section{Anisotropic point jets and isolated poles on fibrations}
\label{sec:anisotropic-growing-jets}

Vanishing to a given order at a point couples the variables through the total degree, even when the metric weight splits as a sum. On a compact product of curves, this coupling becomes a half-space condition on independent monomial indices, and the transition on each curve factor determines the local law in higher dimensions. An isolated pole on a submersion with connected fibers gives a further application: its multiplier ideal does not descend to the base, although the condition it imposes on pulled-back sections does.

Let $q\geq2$. For $1\leq j\leq q$, choose integers $r_j\geq2$, at least two of them unequal, and numbers $\varepsilon_j>0$. Put
\[
 Y=(\mathbb P^1)^q,\qquad x=(0,\ldots,0),\qquad
 d_j=r_j+1,\qquad L=\boxtimes_j\mathcal O_{\mathbb P^1}(d_j).
\]
In the affine coordinates $z_j$ and the product frame $e$, put $t_j=|z_j|^2$ and define the smooth metric and volume by
\[
 \begin{gathered}
 |e|_h^2=\mathrm e^{-\phi},\qquad
 \phi=\sum_j\psi_j(t_j),\qquad
 \psi_j(t)=\log(1+t^{r_j}+\varepsilon_jt^{r_j+1}),\\
 \nu_Y=a(t)\prod_j\mathrm dA_j,\qquad
 a(t)=\prod_j a_j(t_j),\qquad a_j(t)=(1+t)^{-2},
 \end{gathered}
\]
where $\mathrm dA_j$ is Euclidean area in $z_j$. In particular $\nu_Y(Y)=\pi^q$ and $a(0)=1$. Define
\[
 H_j(t)=t\psi_j'(t),\qquad G_j(t)=tH_j'(t),\qquad
 g_j(t)=G_j(t)/t,\qquad \theta=c_1(L,h).
\]
The metric extends across infinity: there its weight, in terms of $u=|z_j|^{-2}$, is $\log(\varepsilon_j+u+u^{d_j})$, which has positive curvature at infinity. For $t>0$, $H_j(t)$ and $G_j(t)$ are the mean and the variance of the exponents $0,r_j,d_j$ taken with probabilities proportional to $1,t^{r_j},\varepsilon_jt^{d_j}$. Hence $G_j>0$, and $H_j$ increases from $0$ to $d_j$. Near zero,
\begin{equation}\label{eq:anisotropic-weight-expansion}
 \begin{aligned}
 H_j(t)&=r_jt^{r_j}+(r_j+1)\varepsilon_jt^{r_j+1}+O(t^{2r_j}),\\
 G_j(t)&=r_j^2t^{r_j}+(r_j+1)^2\varepsilon_jt^{r_j+1}+O(t^{2r_j}).
 \end{aligned}
\end{equation}
Thus $\theta$ is semipositive, and its rank drops along the coordinate divisors $\{z_j=0\}$.

For integers $k_m\to\infty$ with $k_m=o(m)$, let $B_m$ be the density of
\[
 V_m=H^0(Y,L^m\otimes\mathfrak m_x^{k_m}),\qquad
 \|s\|^2=\int_Y|s|_{h^m}^2\nu_Y.
\]
The ideal $\mathfrak m_x^{k_m}$ of the point $x$ imposes vanishing of total order $k_m$, rather than separate orders in the individual variables. As before, $\Phi$ is the standard Gaussian distribution function.

\Needspace{17\baselineskip}
\begin{theorem}\label{thm:anisotropic-growing-jets}\index{Bergman kernel!anisotropic growing jets}
Let $z=z^{(m)}$ be points of the affine chart, with $t_j=|z_j|^2$, such that $t_j\to0$ and $\mu_j:=mH_j(t_j)$ satisfies
$\mu_j/k_m\to\alpha_j>0$, where $\sum_j\alpha_j=1$. Set $V=\sum_jmG_j(t_j)$. If $(\sum_j\mu_j-k_m)/\sqrt V\to\tau\in\mathbb R$, then
\begin{equation}\label{eq:anisotropic-growing-transition}
 \frac{\pi^q a(t)B_m(z)}{m^q\prod_jg_j(t_j)}\longrightarrow\Phi(\tau).
\end{equation}
The arguments of the coordinates $z_j$ are arbitrary. For all sufficiently large $m$,
\begin{equation}\label{eq:anisotropic-jet-dimension}
 N_m:=\dim V_m=\prod_j(d_jm+1)-\binom{k_m+q-1}{q},
 \qquad N_m/m^q\longrightarrow D:=\prod_jd_j.
\end{equation}
Moreover $m^{-q}B_m\nu_Y\to\theta^q/q!$ in total variation.
\end{theorem}
\begin{proof}
The proof reduces the kernel to a probability for a sum of independent variables, one for each factor, and applies the one-variable transition of \cref{sec:compact-growing-jets} to each of them. In each variable, the affine coefficient of a global section of $L^m$ is entire, and its expression in the frame at infinity must be holomorphic, so it is a polynomial of degree at most $d_jm$. Hence the monomials $z^\ell e^m$, $0\leq\ell_j\leq d_jm$, form a basis of the global sections. A polynomial vanishes to order at least $k_m$ at $x$ exactly when all its monomials have total degree at least $k_m$, so the monomials with $\sum_j\ell_j\geq k_m$ form a basis of $V_m$. Once $k_m\leq\min_j(d_jm)+1$, which holds for large $m$, all monomials of total degree below $k_m$ are global, and counting them gives \eqref{eq:anisotropic-jet-dimension}.

Invariance under the rotation of each coordinate makes these monomials orthogonal. Their squared norms are $\prod_jJ_{j,m,\ell_j}$, where
\[
 J_{j,m,\ell}=\pi\int_0^\infty t^\ell\mathrm e^{-m\psi_j(t)}a_j(t)\,\mathrm dt,
 \qquad
 w_{j,m,\ell}(t)=\frac{t^\ell\mathrm e^{-m\psi_j(t)}}{J_{j,m,\ell}}.
\]
These integrals converge even for $\ell=d_jm$. Write $K_{j,m}=\sum_{\ell=0}^{d_jm}w_{j,m,\ell}$, and let $\mathsf L_1,\ldots,\mathsf L_q$ be independent random variables with $\Pr\{\mathsf L_j=\ell\}=w_{j,m,\ell}(t_j)/K_{j,m}(t_j)$ for $0\leq\ell\leq d_jm$. Then
\begin{equation}\label{eq:anisotropic-kernel-halfspace}
 B_m(z)=\prod_jK_{j,m}(t_j)\,
             \Pr\left\{\sum_j\mathsf L_j\geq k_m\right\}.
\end{equation}

Each factor $(\mathbb P^1,\mathcal O(d_j))$, with the metric $\mathrm e^{-\psi_j(|z_j|^2)}$ and the volume $a_j\,\mathrm dA_j$, satisfies the hypotheses of \cref{thm:curvature-boundary-jets}, and $K_{j,m}$ is its full density. Put $V_j=mG_j(t_j)$. Since $\partial_{z_j}\partial_{\bar z_j}\psi_j(|z_j|^2)=g_j(t_j)$, the full density \eqref{eq:mesoscopic-full-density} gives $K_{j,m}=mg_j(t_j)/(\pi a_j(t_j))(1+o(1))$. For any fixed $y\in\mathbb R$, apply the radial formula \eqref{eq:radial-growing-jet-argument} to the integer order
\[
 b_{j,m}=\lfloor\mu_j+y\sqrt{V_j}\rfloor+1.
\]
This order tends to infinity and is sublinear. The selected kernel of the factor, divided by its full kernel, is $\Pr\{\mathsf L_j\geq b_{j,m}\}$, and it tends to $\Phi(-y)$. Hence
\[
 \frac{\mathsf L_j-\mu_j}{\sqrt{V_j}}\ \Longrightarrow\ N(0,1).
\]
that is, these normalized variables converge in distribution to a standard Gaussian variable. Here $\mu_j$ and $\sqrt{V_j}$ serve only as centering and scale; we do not need the exact mean or variance of $\mathsf L_j$. Since $V_j/V\to r_j\alpha_j/\sum_i r_i\alpha_i>0$, independence gives a standard Gaussian limit for $(\sum_j\mathsf L_j-\sum_j\mu_j)/\sqrt V$. To see this, use the product of the limiting Gaussian laws on bounded rectangles, and then tightness and the fact that the boundary of a half-space has zero Gaussian measure. Consequently the probability in \eqref{eq:anisotropic-kernel-halfspace} tends to $\Phi(\tau)$, which proves the local assertion for every sequence $k_m\to\infty$ with $k_m=o(m)$.

For the global assertion we also need the full kernel at a fixed $t>0$, where the formula for shrinking radii does not apply. We suppress the index of the factor and set $F(u)=\psi(\mathrm e^u)$. Then $F'(u)=H(\mathrm e^u)$ and $F''(u)=G(\mathrm e^u)>0$. For $\ell/m$ in a compact subinterval of $(0,d)$, the saddle $s=H^{-1}(\ell/m)$ stays compact, and Laplace expansion gives uniformly
\[
 J_{m,\ell}=\pi a(s)s\sqrt{\frac{2\pi}{mG(s)}}
       \mathrm e^{\ell\log s-m\psi(s)}(1+o(1)).
\]
Indeed, the phase in the logarithmic variable has a uniformly positive second derivative there, and outside a fixed neighborhood of the saddle point, convexity and the slopes at the endpoints bound the tails exponentially. For $\ell=mH(t)+O(\sqrt m)$, evaluation yields
\[
 w_{m,\ell}(t)=\frac{\sqrt{mG(t)}}{\pi a(t)t\sqrt{2\pi}}
   \mathrm e^{-(\ell-mH(t))^2/(2mG(t))}(1+o(1)).
\]
For indices with saddle points in a compact interval, the loss in the phase, as a function of $\ell$, has second derivative $1/(mG(s))\geq c/m$, which gives summable Gaussian tails. For the remaining indices, fix small numbers $\delta,\sigma>0$, and bound the norm $J_{m,\ell}$ from below by the integral over an interval of length $1/m$ about $t\mathrm e^{-\sigma}$ or $t\mathrm e^{\sigma}$. If $|\ell/m-H(t)|\geq\delta$, the appropriate choice gains at least $m(\delta\sigma-C\sigma^2)\geq cm$ in the exponent. Each such weight $w_{m,\ell}(t)$ is then at most $Cm\mathrm e^{-cm}$, and there are $dm+1$ indices. Gaussian summation therefore gives $K_m(t)/m\to g(t)/(\pi a(t))$.

In the product chart, $\theta^q/q!=\prod_j(g_j(t_j)/\pi)\mathrm dA_j$, and its mass is $D$ because $\int_0^\infty g_j(t)\,\mathrm dt=H_j(\infty)-H_j(0)=d_j$. The normalized full density $m^{-q}\prod_jK_{j,m}$ therefore tends to the density of $\theta^q/q!$ with respect to $\nu_Y$ almost everywhere; the coordinate divisors have zero volume. The total mass of the normalized full kernel measure is $m^{-q}\prod_j(d_jm+1)\to D$. For the nonnegative densities $f_m,f$ of these measures with respect to $\nu_Y$, dominated convergence applied to $\min(f_m,f)$ gives
\[
 \int|f_m-f|\nu_Y
 =\int f_m\nu_Y+\int f\nu_Y-2\int\min(f_m,f)\nu_Y\longrightarrow0.
\]
Finally, the difference between the full and the selected kernel is nonnegative and has integral $\binom{k_m+q-1}{q}=o(m^q)$. This proves the total-variation convergence.
\end{proof}

\subsection{Corrected and homogeneous centers}
The transition in \cref{thm:anisotropic-growing-jets} is centered at the radii determined by the functions $H_j$. Replacing them by the radii of the homogeneous leading terms of the weight introduces a shift, which is negligible exactly in a range of orders determined by the largest exponent $r_j$.

Fix positive numbers $\alpha_j$ with sum one and a real $\tau$, and set
\begin{equation}\label{eq:anisotropic-corrected-radii}
 R=\sum_j\alpha_jr_j,\qquad T_m=k_m+\tau\sqrt{Rk_m},\qquad
 t_j=H_j^{-1}(\alpha_jT_m/m).
\end{equation}
Then $\sum_j\mu_j=T_m$ and $V\sim Rk_m$, so \eqref{eq:anisotropic-growing-transition} holds for every sublinear growing order. The radii $|z_j|=\sqrt{t_j}$ are asymptotic to $(\alpha_jk_m/(r_jm))^{1/(2r_j)}$.

\Needspace{17\baselineskip}
\begin{proposition}\label{prop:anisotropic-homogeneous-center}
Keep $\alpha_j,\tau,R,T_m$ as above, but use instead the homogeneous radii given by $t_j=(\alpha_jT_m/(r_jm))^{1/r_j}$, and put $r_*:=\max_jr_j$. For the metrics of this section, the law
\[
 \frac{\pi^qB_m(z)}{m^q\prod_jr_j^2|z_j|^{2r_j-2}}\longrightarrow\Phi(\tau)
\]
holds for every fixed $\tau$ if and only if $k_m=o(m^{2/(r_*+2)})$. If $k_m/m^{2/(r_*+2)}\to c\in(0,\infty)$, the left side instead tends to
\[
 \Phi\bigl(\tau+C_\alpha c^{(r_*+2)/(2r_*)}\bigr),\qquad
 C_\alpha=\frac1{\sqrt R}\sum_{j:r_j=r_*}
 (r_*+1)\varepsilon_j(\alpha_j/r_*)^{1+1/r_*}.
\]
If that ratio tends to infinity while $k_m=o(m)$, the limit is $1$.
\end{proposition}
\begin{proof}
By \eqref{eq:anisotropic-weight-expansion}, $\sum_j\mu_j=T_m+\Delta_m$ with
\[
 \Delta_m=\sum_j(r_j+1)\varepsilon_jmt_j^{r_j+1}
                    +O\left(\sum_jmt_j^{2r_j}\right),\qquad V\sim Rk_m.
\]
Since $r_j\geq2$, each error term is negligible compared with the corresponding positive correction. Consequently
\[
 \frac{\Delta_m}{\sqrt{Rk_m}}=(1+o(1))\sum_j
 \frac{(r_j+1)\varepsilon_j(\alpha_j/r_j)^{1+1/r_j}}{\sqrt R}
 k_m^{1/2+1/r_j}m^{-1/r_j}.
\]
At the critical scale only the indices with $r_j=r_*$ contribute in the limit. Below it every term tends to zero, and above it a positive term diverges. The Gaussian limit of the normalized sum proved above applies in all three cases; in the last case its tightness gives probability one. Moreover, $a(t)\to1$ and $g_j(t_j)/(r_j^2t_j^{r_j-1})\to1$, so the two normalizations agree. If $k_m/m^{2/(r_*+2)}$ does not tend to zero, a subsequence has a positive finite or infinite limit, and the law fails at $\tau=0$. This proves both directions of the equivalence for these metrics.
\end{proof}

\subsection{An isolated selecting potential on a submersion}
Let $\pi:X\to Y$ be a surjective holomorphic submersion from a compact complex manifold $X$, with connected fibers of dimension $f\geq1$, and put $N=q+f=\dim X$. Fix a Hermitian form $\omega_X$ on $X$. Use $E=\pi^*L$ with the metric $h_E=\pi^*h$, and a smooth positive volume form $\nu_X$ satisfying
\begin{equation}\label{eq:anisotropic-matched-volume}
 \pi_*\nu_X=\nu_Y.
\end{equation}
Choose $p\in\pi^{-1}(x)$ and submersion coordinates $(z,w)$ centered at $p$, in which $\pi(z,w)=z$. Let $\chi_p$ be a function equal to $\log(|z|^2+|w|^2)$ near $p$, smooth on $X\setminus\{p\}$, and zero outside a larger coordinate ball. Define the selecting potentials, which depend on $m$,
\[
 \varphi_m=\frac{k_m+N-1}{m}\chi_p,\qquad
 \mathcal I_m=\mathcal I(m\varphi_m).
\]
As in \cref{def:multiplier-ideal}, $\mathcal I_m$ consists of the germs $F$ for which $|F|^2\mathrm e^{-m\varphi_m}$ is locally integrable. The selected sections are measured with the smooth norm defined by $h_E^m$ and $\nu_X$.

\Needspace{17\baselineskip}
\begin{theorem}\label{thm:anisotropic-fiber-selector}\index{Multiplier ideal!isolated selecting potential}
For all large $m$, $\varphi_m\in\PSH(X,\omega_X)$ and $\mathcal I_m=\mathfrak m_p^{k_m}$. Neither $\varphi_m$ nor $\mathcal I_m$ is pulled back from $Y$, but pullback gives an isometry of Hilbert spaces and an identity of densities
\[
 H^0(X,E^m\otimes\mathcal I_m)=\pi^*V_m,\qquad
 B_{X,m}=B_m\circ\pi.
\]
Thus the local laws above hold for every sequence of points of $X$ whose projections satisfy their hypotheses. Moreover $\log N_m/\log m\to q<\dim X$, and
\begin{equation}\label{eq:anisotropic-lifted-measure}
 m^{-q}B_{X,m}\nu_X\longrightarrow
 \left(\frac{\theta^q/q!}{\nu_Y}\circ\pi\right)\nu_X
 \quad\hbox{in total variation}.
\end{equation}
Normalizing by $N_m$ instead gives the limit divided by $D$, a probability measure.
\end{theorem}
\begin{proof}
The function $\chi_p$ satisfies $\ddc\chi_p\geq-C\omega_X$, because it is plurisubharmonic near its pole and smooth elsewhere. Since its coefficient $(k_m+N-1)/m$ tends to zero, $\varphi_m$ is $\omega_X$-plurisubharmonic for large $m$. A germ of order $\ell$ at $p$ has a nonzero initial homogeneous term, so the integral of its squared modulus over the sphere of radius $s$ is $cs^{2\ell}(1+o(1))$ with $c>0$. Local integrability is therefore equivalent to
\[
 \int_0^{s_0}s^{2\ell-2(k_m+N-1)+2N-1}\,\mathrm ds
 =\int_0^{s_0}s^{2\ell-2k_m+1}\,\mathrm ds<\infty,
\]
that is, to $\ell\geq k_m$. Away from $p$ the ideal is the unit ideal. This proves $\mathcal I_m=\mathfrak m_p^{k_m}$. The ideal is proper at $p$ but trivial at the other points of the same fiber, which is impossible for an ideal pulled back from $Y$; likewise, the pole of $\varphi_m$ occurs only at $p$.

By the maximum modulus principle, every section of $\pi^*L^m$, written in a trivialization pulled back from the base, is constant on each compact connected fiber. Local holomorphic sections of the submersion show that the resulting coefficient on the base is holomorphic, and constancy along the fibers makes these coefficients glue. In the submersion coordinates, a pulled-back coefficient depends only on $z$, so its order at $p$ equals its order at $x$. This proves the identity of the selected spaces. Fiber integration and \eqref{eq:anisotropic-matched-volume} give
\[
 \int_X|\pi^*s|_{h_E^m}^2\nu_X=\int_Y|s|_{h^m}^2\nu_Y.
\]
Pulling back an orthonormal basis proves the identity of densities. If $b_m$ and $b$ are the densities, with respect to $\nu_Y$, of the measures in the total-variation limit of \cref{thm:anisotropic-growing-jets}, then the corresponding densities on $X$ are $b_m\circ\pi$ and $b\circ\pi$, and their $L^1(\nu_X)$ distance equals $\int_Y|b_m-b|\nu_Y$. This proves \eqref{eq:anisotropic-lifted-measure}. The dimension formula \eqref{eq:anisotropic-jet-dimension} gives $\log N_m/\log m\to q$ and the final normalization.
\end{proof}

Volume forms satisfying \eqref{eq:anisotropic-matched-volume} exist on every such submersion. For any smooth positive volume form $\mu_X$, write $\pi_*\mu_X=b\nu_Y$. Submersion charts and a finite partition of unity express $b$ as a sum of integrals of smooth compactly supported densities in the fiber variables, and differentiation under the integral sign shows that $b$ is smooth. Since the fibers are nonempty, $b>0$. Thus $\nu_X=\mu_X/(b\circ\pi)$ satisfies \eqref{eq:anisotropic-matched-volume}. For a product $X=Y\times F$, the fiber $F$ may be any compact connected complex manifold, K\"ahler or not. On a product, both the volumes $\nu_Y\otimes\nu_F$ with $\nu_F(F)=1$ and the nonproduct volumes $W(y,u)\nu_Y\nu_F$, with $W>0$ and $\int_FW(y,u)\nu_F=1$, satisfy \eqref{eq:anisotropic-matched-volume}.

For an arbitrary smooth positive volume $\nu'_X$, the pulled-back norm becomes the norm on the base for the volume $\pi_*\nu'_X$, whereas the local laws above were proved for the specific volume $\nu_Y$. Since the selecting potentials vary with $m$, \cref{thm:anisotropic-fiber-selector} does not settle \cref{qu:fiber-selected-approximation} for fixed singularities. Weights with mixed or angular remainder terms, observation points with some $\alpha_j=0$, and other ideals on the base are not treated here. The lifted measure keeps the chosen conditional fiber volume; it is not a Monge--Amp\`ere measure in the dimension of $X$.
